% Financial Sentiment Analyzer — tool overview (Beamer)
% Build: pdflatex tool-overview.tex  (run twice for TOC)
\documentclass[aspectratio=169,11pt]{beamer}

\usepackage[T1]{fontenc}
\usepackage[utf8]{inputenc}
\usepackage{lmodern}
\usepackage{microtype}
\usepackage{booktabs}
\usepackage{hyperref}
\usepackage{listings}
\usepackage{xcolor}
\usepackage{graphicx}

\usetheme{Madrid}
\usecolortheme{default}

\title{Financial Sentiment Analyzer}
\subtitle{Architecture, pipeline, and backtest evaluation}
\author{FYP / internal technical overview}
\date{\today}

\definecolor{codebg}{RGB}{248,248,255}
\lstset{
  basicstyle=\ttfamily\footnotesize,
  backgroundcolor=\color{codebg},
  frame=single,
  framerule=0.4pt,
  rulecolor=\color{gray!60},
  breaklines=true,
  breakatwhitespace=true,
  columns=fullflexible,
  keepspaces=true,
  showstringspaces=false,
  tabsize=2,
  keywordstyle=\color{blue!70!black},
  commentstyle=\color{gray!70},
  stringstyle=\color{teal!60!black},
}

\hypersetup{colorlinks=true,linkcolor=blue,urlcolor=blue}

\begin{document}

\begin{frame}
  \titlepage
\end{frame}

\begin{frame}{Outline}
  \tableofcontents
\end{frame}

%-------------------------------------------------------------------------
\section{Purpose and capabilities}

\begin{frame}{What this tool does}
  \begin{itemize}
    \item \textbf{Ingests} financial news and links articles to \textbf{companies} (tickers).
    \item \textbf{Scores sentiment} with a \textbf{FinBERT}-style transformer (base + optional \textbf{fine-tuned} weights).
    \item \textbf{Stores} structured results in \textbf{PostgreSQL} (articles, scores, OHLCV, correlations).
    \item \textbf{Serves} a \textbf{FastAPI} backend: trends, company sentiment, correlation metrics, optional social sentiment.
    \item \textbf{Explains} context via \textbf{RAG chat} (embeddings + \textbf{Groq} LLM).
    \item \textbf{Evaluates} the sentiment signal in a \textbf{long-horizon trading simulation} (LLM agents; treatment vs.\ control).
  \end{itemize}
\end{frame}

\begin{frame}{Who it is for}
  \begin{itemize}
    \item \textbf{End users} (via frontend): explore news sentiment, trends, and chat with retrieved evidence.
    \item \textbf{You (creator)}: offline artefacts --- fine-tuning jobs, \textbf{simulation reports} under \texttt{data/simulations/}, not exposed as a public product surface.
  \end{itemize}
\end{frame}

%-------------------------------------------------------------------------
\section{System architecture}

\begin{frame}{High-level architecture}
  \begin{center}
    \small
    \begin{tabular}{@{}l l@{}}
      \toprule
      \textbf{Layer} & \textbf{Components} \\
      \midrule
      Ingestion & Celery workers: news collection, market data, backfills \\
      NLP & \texttt{SentimentAnalyzer} (FinBERT / active fine-tuned checkpoint) \\
      Storage & PostgreSQL (articles, \texttt{sentiment\_results}, \texttt{market\_data}, \ldots) \\
      Search / chat & Chroma + sentence embeddings; Groq for generation \\
      API & FastAPI (\texttt{/api/news}, \texttt{/api/chat}, \texttt{/api/correlations}, \ldots) \\
      Evaluation & Trader-agent \textbf{backtest} (optional Fireworks / Groq for decisions) \\
      \bottomrule
    \end{tabular}
  \end{center}
\end{frame}

\begin{frame}[fragile]{API surface (excerpt)}
  Routers are mounted in \texttt{app/main.py}:
\begin{lstlisting}[language=Python]
app.include_router(news.router, prefix="/api/news")
app.include_router(chat.router, prefix="/api/chat")
app.include_router(finetuning.router, prefix="/api/finetuning")
app.include_router(simulations.router, prefix="/api/simulations")
\end{lstlisting}
  \vfill
  \footnotesize Health: \texttt{GET /api/health} \quad Status (model names): \texttt{GET /api/status}
\end{frame}

%-------------------------------------------------------------------------
\section{Sentiment pipeline}

\begin{frame}{From raw article to score}
  \begin{enumerate}
    \item News rows live in \texttt{news\_articles} (title, content, \texttt{publication\_date}, \ldots).
    \item \texttt{analyze\_sentiment\_task} batches texts through the active model.
    \item Results persist to \texttt{sentiment\_results} (\texttt{sentiment\_label}, softmax scores, confidence).
    \item Company linkage uses \texttt{article\_companies} so tickers can aggregate sentiment over time.
  \end{enumerate}
\end{frame}

\begin{frame}[fragile]{Sentiment task (core idea)}
\begin{lstlisting}[language=Python]
# app/workers/tasks.py - batched GPU inference, idempotent writes
texts = [f"{a.title}. {a.content}" if a.content else a.title
         for a in articles]
predictions = analyzer.batch_analyze(texts, batch_size=32)
# ... map each prediction -> SentimentResult(article_id=..., ...)
session.commit()
\end{lstlisting}
\end{frame}

\begin{frame}{Fine-tuning (optional)}
  \begin{itemize}
    \item \textbf{Base model} configurable (e.g.\ FinBERT tone).
    \item \textbf{Fine-tuning jobs} train on a parquet dataset; best checkpoint path stored on disk.
    \item \textbf{Active model}: whichever job is marked active drives inference for new batches.
    \item Keeps \textbf{label mapping} consistent between training and inference.
  \end{itemize}
\end{frame}

%-------------------------------------------------------------------------
\section{Simulation and evaluation}

\begin{frame}{Why a trading simulation?}
  \begin{itemize}
    \item Measures whether \textbf{treatment} agents (with sentiment briefings) behave differently from \textbf{control} agents (prices/vol only).
    \item \textbf{No look-ahead}: briefings use information available as of prior close; orders fill at the \textbf{open}.
    \item \textbf{Long transaction}: results commit at end; \textbf{progress snapshots} every N seconds to disk.
  \end{itemize}
\end{frame}

\begin{frame}[fragile]{LLM provider for agents (excerpt)}
\begin{lstlisting}[language=Python]
# app/services/simulation/llm_provider.py
if provider == "fireworks":
    return ChatFireworks(
        model=settings.FIREWORKS_MODEL,
        api_key=settings.FIREWORKS_API_KEY,
        request_timeout=timeout,
        max_retries=max_retries,
    )
\end{lstlisting}
  Structured decisions use \texttt{with\_structured\_output(OrderList, method="json\_schema")} on Fireworks.
\end{frame}

\begin{frame}[fragile]{Structured trade orders}
\begin{lstlisting}[language=Python]
# app/services/simulation/agent.py - Pydantic-enforced output
class OrderList(BaseModel):
    orders: list[StructuredOrder] = Field(default_factory=list)

# LLM returns buy/sell, size_pct in [0,1], reasoning; engine applies clamps.
\end{lstlisting}
\end{frame}

\begin{frame}{Progress artefacts}
  \begin{itemize}
    \item Directory: \texttt{data/simulations/<run\_id>/progress/}
    \item Files: \texttt{snapshot\_NNNN.md/json}, \texttt{latest.md/json}
    \item Cadence: \texttt{SIMULATION\_PROGRESS\_INTERVAL\_SEC} (e.g.\ 600\,s)
  \end{itemize}
\end{frame}

%-------------------------------------------------------------------------
\section{Initial backtest snapshot (run 4)}

\begin{frame}{Snapshot context}
  \footnotesize
  Source: \texttt{data/simulations/4/progress/latest.md} (progress snapshot \#5). \\
  \textbf{Note:} early in the backtest ($\sim$2\% of trading days); rankings are \emph{not} final conclusions.
  \vfill
  \begin{itemize}
    \item Wall-clock (engine): \textbf{57m 20s} \quad ETA (reported): $\sim$\textbf{43h 30m}
    \item Progress: \textbf{day 27 / 1256} (2.1\%) \quad Simulated date: \textbf{2021-06-01}
  \end{itemize}
\end{frame}

\begin{frame}{Leaderboard excerpt (\texttt{latest.md})}
  \footnotesize
  \begin{tabular}{@{}r l l r r@{}}
    \toprule
    Rank & Profile & Variant & Equity & Return \\
    \midrule
    1 & value\_investor & treatment & \$1{,}038.54 & +3.85\% \\
    2 & contrarian & treatment & \$1{,}022.91 & +2.29\% \\
    3 & momentum\_investor & control & \$1{,}015.86 & +1.59\% \\
    4 & contrarian & control & \$1{,}006.56 & +0.66\% \\
    \ldots & \ldots & \ldots & \ldots & \ldots \\
    12 & swing\_trader & control & \$959.61 & $-$4.04\% \\
    \bottomrule
  \end{tabular}
\end{frame}

\begin{frame}{Head-to-head: treatment $-$ control (early)}
  \footnotesize
  \begin{tabular}{@{}l r r r@{}}
    \toprule
    Profile & Treat.\ \% & Ctrl.\ \% & Lift \\
    \midrule
    value\_investor & +3.85 & +0.00 & +3.85 \\
    contrarian & +2.29 & +0.66 & +1.64 \\
    swing\_trader & $-$2.59 & $-$4.04 & +1.45 \\
    news\_driven & $-$3.81 & $-$1.74 & $-$2.07 \\
    momentum\_investor & +0.00 & +1.59 & $-$1.59 \\
    day\_trader & $-$2.08 & $-$1.59 & $-$0.49 \\
    \bottomrule
  \end{tabular}

  \vfill
  \footnotesize Interpret with care: different \textbf{cadences} (daily vs.\ weekly agents), LLM noise, and fees/slippage model are simplified.
\end{frame}

%-------------------------------------------------------------------------
\section{Operations and stack}

\begin{frame}{Typical stack (Docker Compose)}
  \begin{itemize}
    \item \textbf{backend} --- FastAPI + FinBERT load
    \item \textbf{celery-worker} --- sentiment, news, simulation tasks
    \item \textbf{postgres}, \textbf{redis}, \textbf{chromadb}
    \item Volumes: models, datasets, \texttt{./data/simulations} bind-mount for reports
  \end{itemize}
\end{frame}

\begin{frame}{Takeaways}
  \begin{itemize}
    \item End-to-end: \textbf{news} $\rightarrow$ \textbf{FinBERT scores} $\rightarrow$ \textbf{APIs \& RAG} $\rightarrow$ \textbf{optional fine-tune}.
    \item Internal \textbf{simulation} stress-tests whether sentiment briefings affect LLM trading vs.\ controls.
    \item Early run~4 snapshot shows \textbf{mixed lift by profile}; wait for full \texttt{run\_summary.md}.
  \end{itemize}
\end{frame}

\begin{frame}
  \centering
  \vfill
  {\Huge Questions?}
  \vfill
\end{frame}

\end{document}
